\documentclass[10pt]{article}
\usepackage[T1]{fontenc}
\usepackage[utf8]{inputenc}
\usepackage[french]{babel}
\usepackage{amsmath,amssymb,mathrsfs,mathpazo}
\usepackage{geometry}
\geometry{a4paper,margin=2cm}
\usepackage{graphicx,array,tabularx,tabularray,colortbl,multirow,multicol,booktabs,enumitem,pifont,ulem,fancybox,diagbox,cancel}
\usepackage{pstricks,pstricks-add,pst-plot,pst-tree,pst-node,pst-text,pst-func,pst-math,pst-bspline,pst-3dplot,pst-eucl,pst-solides3d}
\let\pscalculate\undefined
\usepackage{pst-bezier}
\usepackage{tikz,pgfplots,tkz-tab}
\usetikzlibrary{arrows,arrows.meta,patterns,calc,positioning,trees,decorations.markings,shapes.geometric,shapes.misc,shapes.arrows,intersections,decorations.pathmorphing,decorations.pathreplacing,backgrounds,shadows,babel}
\pgfplotsset{compat=1.18}
\usepgfplotslibrary{fillbetween}
\usepackage[np]{numprint}
\usepackage{listings,makecell,pas-tableur,twoopt}
\let\barre\undefined
\usepackage[pstricks]{tablvar}
\definecolor{darkgreen}{RGB}{0,100,0}
\providecommand{\red}{\color{red}}
\providecommand{\blue}{\color{blue}}
\providecommand{\green}{\color{green}}
\providecommand{\black}{\color{black}}
\providecommand{\white}{\color{white}}
\providecommand{\gray}{\color{gray}}
\providecommand{\cyan}{\color{cyan}}
\providecommand{\magenta}{\color{magenta}}
\providecommand{\yellow}{\color{yellow}}
\providecommand{\violet}{\color{violet}}
\providecommand{\orange}{\color{orange}}
\providecommand{\euro}{\texteuro}
\providecommand{\EUR}{\texteuro}
\providecommand{\ieme}{\textsuperscript{e}}
\providecommand{\degres}{\ensuremath{^\circ}}
\providecommand{\degree}{\ensuremath{^\circ}}
\providecommand{\vv}[1]{\overrightarrow{#1}}
\providecommand{\Coord}[1]{#1}
\providecommand{\decosix}{}
\providecommand{\starredbullet}{\ensuremath{\bullet}}
\pagestyle{empty}
\setlength{\parindent}{0pt}
\usepackage{adjustbox}\DeclareUnicodeCharacter{2236}{:}\DeclareUnicodeCharacter{2190}{\ensuremath{\leftarrow}}
\providecommand{\R}{\mathbb{R}}
\providecommand{\N}{\mathbb{N}}
\providecommand{\Z}{\mathbb{Z}}
\providecommand{\Q}{\mathbb{Q}}
\providecommand{\D}{\mathbb{D}}
\providecommand{\C}{\mathbb{C}}
\providecommand{\e}{{\mathrm{e}}}
\providecommand{\vect}[1]{\overrightarrow{#1}}
\providecommand{\vectt}[1]{\overrightarrow{\mathrm{#1}}}
\providecommand{\vvt}[1]{\overrightarrow{\mathrm{#1}}}
\providecommand{\np}[2][]{{#2~#1}}
\providecommand{\vv}[1]{{\overrightarrow{#1}}}
\renewcommand{\d}{{\mathrm{d}}}
\providecommand{\degres}{{^\circ}}
\providecommand{\degre}{{^\circ}}
\providecommand{\blue}{\color{blue}}
\providecommand{\red}{\color{red}}
\providecommand{\enskip}{\quad}
\providecommand{\llbracket}{\lbrack\!\lbrack}
\providecommand{\rrbracket}{\rbrack\!\rbrack}
\providecommand{\euro}{{\text{€}}}
\providecommand{\no}{{\text{n°}}}
\providecommand{\interval}[2]{{[#1;#2]}}
\providecommand{\intervFO}[2]{{[#1;#2[}}
\providecommand{\dotfill}{\ldots}

\providecommand{\ds}{\displaystyle}
\providecommand{\textding}[1]{#1}
\providecommand{\Oij}{\(\mathrm{O};\vec{\imath},\vec{\jmath}\)}
\providecommand{\Oijk}{\(\mathrm{O};\vec{\imath},\vec{\jmath},\vec{k}\)}
\providecommand{\Ouv}{\(\mathrm{O};\vec{u},\vec{v}\)}
\providecommand{\intervFF}[2]{\ensuremath{[#1;#2]}}
\renewcommand{\euro}{\ifmmode\text{\texteuro}\else\texteuro\fi}
\DeclareUnicodeCharacter{00B0}{\ifmmode^{\circ}\else\textdegree\fi}
\DeclareUnicodeCharacter{20AC}{\ifmmode\text{\texteuro}\else\texteuro\fi}
% Énoncés : Éducation nationale. Transcription et corrigés : APMEP, Denis Vergès et contributeurs. Découpage et classement : CoopMaths, Sébastien Lozano, Éric Elter et contributeurs. Adaptation de la présentation : Maths Explorer.
% CC BY-SA : https://coopmaths.fr/www/about/
% https://coopmaths.fr/alea/static/bac/2025/tex/bac_2025_06_sujet2_etrangers_1.tex
\begin{document}
\noindent\textit{APMEP / CoopMaths. CC BY-SA. Présentation adaptée par Maths Explorer.}\par\medskip
On se propose de comparer l'évolution d'une population animale dans deux milieux distincts A et B.

Au 1\up{er} janvier 2025, on introduit \np{6000} individus dans chacun des milieux A et B.

\medskip

\textbf{Partie A}

\medskip

Dans cette partie, on étudie l'évolution de la population dans le milieu A.

On suppose que dans ce milieu, l'évolution de la population est modélisée par une suite géométrique $(u_n)$ de premier terme $u_0 = 6$ et de raison $0,93$.

Pour tout entier naturel $n,\: u_n$ représente la population au 1\up{er} janvier de l'année $2025+n$, exprimée en millier d'individus.

\medskip

\begin{enumerate}
\item Donner, selon ce modèle, la population au 1\up{er} janvier 2026.
\item Pour tout entier naturel $n$, exprimer $u_n$ en fonction de $n$.
\item Déterminer la limite de la suite $(u_n)$.

Interpréter ce résultat dans le contexte de l'exercice.
\end{enumerate}

\medskip

\textbf{Partie B}

\medskip

Dans cette partie, on étudie l'évolution de la population dans le milieu B.

On suppose que dans ce milieu, l'évolution de la population est modélisée par la suite $(v_n)$ définie par 

\begin{center}$v_0 = 6$ et pour tout entier naturel $n,\: v_{n+1} = -0,05v_n^2 +1,1v_n$.\end{center}

Pour tout entier naturel $n, v_n$ représente la population au 1\up{er} janvier de l'année $2025 + n$, exprimée en millier d'individus.

\medskip

\begin{enumerate}
\item Donner, selon ce modèle, la population au 1\up{er} janvier 2026.
\end{enumerate}

Soit $f$ la fonction définie sur l'intervalle $[0~;~+\infty[$ par 

\[f(x) = -0,05x^2 + 1,1x.\]

\begin{enumerate}[resume]
\item Démontrer que la fonction $f$ est croissante sur l'intervalle [0~;~11].
\item Démontrer par récurrence que pour tout entier naturel $n$, on a 

\[2 \leqslant v_{n+1} \leqslant v_n \leqslant 6.\]

\item En déduire que la suite $(v_n)$ est convergente vers une limite $\ell$.
\item 
	\begin{enumerate}
		\item Justifier que la limite $\ell$ vérifie $f(\ell) = \ell$ puis en déduire la valeur de $\ell$.
		\item Interpréter ce résultat dans le contexte de l'exercice.
	\end{enumerate}
\end{enumerate}

\medskip

\textbf{Partie C}

Cette partie a pour but de comparer l'évolution de la population dans les deux milieux.

\begin{enumerate}
\item En résolvant une inéquation, déterminer l'année à partir de laquelle la population du milieu A sera strictement inférieure à \np{3000} individus.
\item À l'aide de la calculatrice, déterminer l'année à partir de laquelle la population du milieu B sera strictement inférieure à \np{3000} individus.
\end{enumerate}
\begin{minipage}{0.78\linewidth}
\begin{enumerate}[resume]
\item Justifier qu'à partir d'une certaine année, la population du milieu B dépassera la population du milieu A.
\item On considère le programme Python ci-contre.
	\begin{enumerate}
		\item Recopier et compléter ce programme afin qu'après exécution, il affiche l'année à partir de laquelle la population du milieu B est strictement supérieure à la population du milieu A.
		\item Déterminer l'année affichée après exécution du programme.
	\end{enumerate}
\end{enumerate}
\end{minipage}\hfill
\begin{minipage}{0.19\linewidth}
\begin{tabular}{|l|}\hline
n=0 \\
u = 6\\
v = 6 \\
\textbf{while} \ldots :\\
\qquad u = \ldots\\
\qquad v= \ldots\\
\qquad n = n+1\\
\textbf{print} (2025 + n)\\ \hline
\end{tabular}
\end{minipage}

\end{document}
